% Section 6: Reflection, shift and the Toda structure (from research/G1/RH/spectral/JZC_M4.md Sec. 7-8 and JZC_SIGMA.md).
\section{Reflection, shift and the Toda structure}\label{sec:shiftsec}

Theorem~\ref{thm:N6} computes the layers of the lattice expansion to fixed orders and observes that they are even in $1/N$.
This section explains the evenness and the natural index to all orders. The exponent $\alpha$ plays two roles in the
residue recursion, and separating them gives a two-parameter family of formal problems (\S\ref{sec:decoupled}). On this
family there are two exact symmetries: a reflection, which is algebra (Theorem~\ref{thm:reflect}), and a shift, which we
prove by steepest descent of the same lattice problem in a second, non-integer frame (Theorem~\ref{thm:shift}). Together they
give the evenness of the integer layer to all orders (Corollary~\ref{cor:even}). The zeta layers are Toda flows of the integer
layer (\S\ref{sec:toda}).

\subsection{The decoupled model}\label{sec:decoupled}
The exponent $\alpha$ enters the jump data of Lemma~\ref{lem:N6a} in two ways: as the Bessel order $\nu$ in the coefficients
$a_k(\nu),b_k(\nu)$, which are polynomials in $\nu^2$, and as the Szeg\H{o} exponent $\beta$ in $\rho^2=\phi^{2(\beta+1)}$ at the hard
edge and in the factor $\phi^{-2(\beta+1)}$ at the soft edge, where $\phi(z)=(z^{1/2}+(z-4c)^{1/2})/(2c^{1/2})$ is the exterior
conformal map of $[0,4c]$. For the lattice $\nu=\beta=\alpha$. Keeping $\nu$ and $\beta$ independent, the residue
recursion~\eqref{eq:recursion} and the extraction of \S\ref{sec:recursion}, with the gauge of Lemma~\ref{lem:gauge} for the
band-factor parameters $\mu$, define formal series
\begin{equation}\label{eq:decoupled}
  \frac{n^p\,b(n;\beta,\nu;\mu)}{c}=\sum_{\br}\mu^{\br}\,\cD_{\br}(n;\beta,\nu),\qquad
  B(n;\beta,\nu):=\cD_0(n;\beta,\nu)=\sum_{k\ge0}d_k(p,\beta,\nu)\,n^{-k},
\end{equation}
whose truncations at order $n^{-K}$ are the polynomials $\cD_{\br,K}$ of \S\ref{sec:recursion} when $\beta=\nu=\alpha$. We call
this the \emph{decoupled model} $M(\beta,\nu)$ and $B$ its \emph{integer layer}. Each $d_k$ is a rational function of
$(p,\beta,\nu)$, $d_0=1$: with $c=1$ the recursion and the extraction are field operations on Laurent data that are rational
in $p$ (Lemma~\ref{lem:N6d}), polynomial in $\nu$ through $a_k(\nu),b_k(\nu)$, and polynomial in $\beta$, since $\beta$ enters
only through binomial series of $(1+\text{series without constant term})^{2(\beta+1)}$.

\begin{theorem}[reflection]\label{thm:reflect}
As formal series, $b(n;\beta,\nu;0)=b(-n;-\beta-2,\nu;0)$; that is, $d_k(p,-\beta-2,\nu)=(-1)^kd_k(p,\beta,\nu)$ for all $k$.
\end{theorem}

\begin{proof}
With $T=\frac1{\sqrt2}\left(\begin{smallmatrix}1&1\\ i&-i\end{smallmatrix}\right)$ we have $\sigma_2T=T\sigma_3$, so
$P_0^\infty=T\mathfrak a_\infty^{\sigma_3}T^{-1}$ and $P_1=-c\sigma_2$ commute with $\sigma_2$. The core matrices $X_k$ and
$\sigma_3\Xi_k\sigma_3$ lie in the span of $I,\sigma_2$ for even $k$ and of $\sigma_3,i\sigma_1$ for odd $k$, so conjugation by
$\sigma_2$ multiplies them by $(-1)^k$, while $\sigma_2d^{\sigma_3}\sigma_2=d^{-\sigma_3}$ for scalars $d$. Hence
$\sigma_2\Delta_k(\beta)n^{-k}\sigma_2=\Delta_k(-\beta-2)(-n)^{-k}$, the recursion is covariant under conjugation by a constant
matrix, and $(Q_1)_{12}(Q_1)_{21}$ is invariant under conjugation by $\sigma_2$.
\end{proof}

\begin{lemma}[self-similar Buyarov--Rakhmanov identity]\label{lem:BR}
Let $\Phi$ denote the phase of the lattice equilibrium in the scaled frame, with hard-edge branch $2f^{1/2}$ and soft-edge
branch $\tilde\varphi$. Then $(1+pz\partial_z)\Phi=-2\log\phi$ up to the additive constants of the branches; in particular
$(1+pz\partial_z)(2f^{1/2})=2\operatorname{arsinh}(v/2)$ at $z=-v^2$ ($c=1$).
\end{lemma}

\begin{proof}
At the hard edge $2f^{1/2}=\pi\sqrt{-z}\,\mathfrak a(z)$, and the series of $\mathfrak a$ has the factors $1/(a'+k)$ with
$a'=\frac1p+\frac12$. Since $(1+pz\partial_z)[\sqrt{-z}\,z^k]=p(a'+k)\sqrt{-z}\,z^k$, these factors cancel term by term and
leave the arcsine series. The soft-edge branch is the analytic continuation of the same phase.
\end{proof}

Under the self-similar scaling of the lattice problem, $1+pz\partial_z$ is the derivative in $n$ of $n\Phi(xn^p)$; so
Lemma~\ref{lem:BR} says that the $n$-flow of the phase is generated by $\log\phi$, which also generates the $\beta$-flow
($\partial_\beta\rho^2=2\log\phi\cdot\rho^2$). This suggested the following theorem; its proof does not use the lemma.

\subsection{The shift theorem}\label{sec:shift-thm}
\begin{theorem}[shift]\label{thm:shift}
Identically in $(p,\beta,\nu)$, as formal series:
\begin{enumerate}
\item[(i)] $B(n;\beta,\nu)=(n/N)^pF_p(N;\nu)$ with $N:=n+\frac{\beta+1}2$, where $F_p(N;\nu)$ is a formal power series in $1/N$
that does not depend on $\beta$; equivalently, for every real $\gamma$ (the \emph{frame identity})
\begin{equation}\label{eq:FI}
  B(n+\gamma;\beta-2\gamma,\nu)=\Bigl(1+\frac\gamma n\Bigr)^pB(n;\beta,\nu),
\end{equation}
the left side re-expanded in $1/n$;
\item[(ii)] for $\br\in\{e_1,e_2,2e_1\}$ the same holds with the exponent $p$ replaced by $(1+|\br|_w)p$: the series
$(N/n)^{(1+|\br|_w)p}\cD_{\br}(n;\beta,\nu)$, re-expanded in $1/N$, does not depend on $\beta$.
\end{enumerate}
\end{theorem}

The infinitesimal form of~\eqref{eq:FI} is the flow identity $\partial_nB-2\partial_\beta B=(p/n)B$. The equivalence of~(i)
and~\eqref{eq:FI} is algebra: if (i) holds, then $(n+\gamma)+\frac{\beta-2\gamma+1}2=N$ and
$B(n+\gamma;\beta-2\gamma,\nu)=((n+\gamma)/N)^pF_p(N)=(1+\gamma/n)^p(n/N)^pF_p(N)$. Conversely, given~\eqref{eq:FI}, fix
$\beta_0$ and put $\gamma:=(\beta_0-\beta)/2$; then $(N/n)^pB(n;\beta,\nu)=(N/(n-\gamma))^pB(n-\gamma;\beta_0,\nu)$, and
$n-\gamma=N-\frac{\beta_0+1}2$ does not involve $\beta$. All re-expansions are between formal series in $1/n$, $1/(n-\gamma)$,
$1/N$, which differ by constants.

\begin{corollary}[evenness]\label{cor:even}
$F_p(N;\nu)=F_p(-N;\nu)$, and $F_p$ depends on $\nu$ only through $\nu^2$. Hence $F_p(N;\nu)=1+\sum_{k\ge1}f_{2k}N^{-2k}$ with
each $f_{2k}$ rational in $p$ and polynomial in $\nu^2$. For the lattice and every $p\ne1$, the integer layer of
Theorem~\ref{thm:N6} is, for every $K$,
\[
  \cD_{0,K}(n)=\Bigl(\frac nN\Bigr)^p\Bigl[1+\sum_{2k\le K}\frac{f_{2k}(p,\alpha^2)}{N^{2k}}\Bigr]+O(N^{-K-1}),\qquad N=n+\frac\alpha2+\frac12 ;
\]
in particular all odd coefficients $f_{2k+1}$ vanish, and the natural index of the lattice is the Szeg\H{o} shift $\frac{\beta+1}2$ at
$\beta=\alpha$.
\end{corollary}

\begin{proof}
By Theorem~\ref{thm:reflect}, $B(n;\beta,\nu)=B(-n;-\beta-2,\nu)$. The natural index of the right side is
$-n+\frac{-\beta-1}2=-N$, so Theorem~\ref{thm:shift}(i) gives $(n/N)^pF_p(N)=((-n)/(-N))^pF_p(-N)$. The coefficients $a_k(\nu)$,
$b_k(\nu)$ depend on $\nu$ only through $\nu^2$, and $\beta$ no longer occurs in $F_p$. Each $d_k$ is polynomial in $\nu$, hence
so is each $f_{2k}$.
\end{proof}

Corollary~\ref{cor:even} settles the integer-layer part of Remark~\ref{rem:N6}(a)--(c): $f_7=0$ and the evenness in $\alpha$ hold to
all orders, while the values $f_8,f_{10},\dots$ and the observed degree $k$ of $f_{2k}$ in $\alpha^2$ remain numerical. The $\alpha$-even
part has the heuristic explanation of Remark~\ref{rem:N6}(c): $\alpha$ enters the integer layer through $\nu^2$ and through $\beta$,
and the shift absorbs $\beta$.

\subsection{Proof of the shift theorem: a second frame}\label{sec:shift-proof}
The decoupled model is not related to itself by any transformation of the formal recursion that we know. In particular,
the Christoffel shift $x\sigma_{p,s}=\sigma_{p,s+p}$ of \cite{PaperII} moves the Bessel order with the Szeg\H{o} exponent,
$(\beta,\nu)\mapsto(\beta+1,\nu+1)$, and any polynomial modification of the measure that is regular and nonzero at $0$ moves
neither. The proof instead analyses the same exact object twice: the degree-$n$ Riemann--Hilbert problem $Y$ of the lattice
is steepest-descended in the frame of an arbitrary real index $m=n+\gamma$, and the two resulting expansions of the same
number $b_n$ are compared.

Let $U_m(z):=m^{pn\sigma_3}Y(zm^{-p})$, the rescaling of \cite{PaperII} with $m^p$ in place of $n^p$; the degree stays $n$. Put
$t_m:=mz^{-1/p}$, so that the nodes $X_k=(m/k)^p$ satisfy $t_m(X_k)=k$, and let $g$, $\phi$, $\tilde\varphi$ be the lattice
$g$-function and phases of \cite{PaperII}, which do not depend on the index. With the principal branch of $z^\gamma$ (cut
$(-\infty,0]$) and the constant $C=m^{p\gamma\sigma_3}$, consider
\[
  \tilde U:=C\,U_mz^{\gamma\sigma_3}C^{-1},\qquad \tilde Uz^{-m\sigma_3}=C(U_mz^{-n\sigma_3})C^{-1}\to I\quad(z\to\infty).
\]
The conjugation by a constant diagonal matrix leaves the product of the off-diagonal entries of the $1/z$ coefficient unchanged,
so $\tilde U$ still encodes $m^{2p}b_n^2$.

\begin{lemma}[the $m$-frame]\label{lem:frame}
(i) The residue conditions of $\tilde U$ are those of the lattice measure $\sigma_{p,\tilde s}$, $\tilde s=s-2p\gamma$, in the frame
of index $m$, so that its band weight has the exponent $\tilde\alpha:=\alpha-2\gamma$. After Steps~1--4 of \cite[Section~5]{PaperII} with
$n$ replaced by $m$ in every place except the degree, the jumps are those of \cite{PaperII} with weight exponent $\tilde\alpha$, there
is no jump on $(-\infty,0)$, and the $(2,1)$ entry of the lip jumps in the hard-edge disc is $\omega^\natural e^{-4\zeta^{1/2}}$,
$\zeta=m^2f$, with $\omega^\natural=\omega^{(\tilde\alpha)}e^{\mp2\pi i\gamma}$ in $\C_\pm$, whose boundary values on the band are
$e^{\mp i\pi\alpha}$. That is, the $m$-frame problem has Bessel order $\alpha$ and Szeg\H{o} exponent $\alpha-2\gamma$.

(ii) Let $p>1$, $|\alpha|<1$ and $0\le\gamma<(1-|\alpha|)/2$. Then for all $K,J,B\ge1$, with $\mu_{m,j}=2\zeta(-jp)/(jm^{jp})$,
\[
  \frac{m^pb_n}c=\sum_{|\br|\le B}\mu_m^{\br}\,\cD_{\br,K}(m;\alpha-2\gamma,\alpha)+O\bigl(m^{-K-1}+m^{-2-(B+1)p}+m^{-1-(J+1)p}\bigr).
\]
\end{lemma}

\begin{proof}[Proof sketch]
(i) The residue of $\tilde U_{j2}$ at $X_k$ acquires the factor $(c_1/c_2)X_k^{-2\gamma}=k^{2p\gamma}$, which turns the mass $k^{-s}$ into
$k^{-\tilde s}$. The computation of the jumps in \cite[Section~5]{PaperII} uses the integrality of the index only on $(-\infty,0)$,
where $e^{m\phi\sigma_3}$ jumps by $e^{-2\pi im\sigma_3}$ and the cut of $z^{\gamma\sigma_3}$ contributes $e^{2\pi i\gamma\sigma_3}$; the
product is $e^{-2\pi in\sigma_3}=I$ because the degree $n$ is an integer. In the hard-edge disc $\phi=\mp i\pi-2f^{1/2}$ in $\C_\pm$
\cite{PaperII}, so $e^{2m\phi}=e^{\mp2\pi i\gamma}e^{-4m f^{1/2}}$; the Szeg\H{o} quotient $\omega^{(\tilde\alpha)}$ tends to
$e^{\mp i\pi\tilde\alpha}$ on the band, and $e^{\mp i\pi(\tilde\alpha+2\gamma)}=e^{\mp i\pi\alpha}$. At the soft edge every jump is written
with $\tilde\varphi$, which does not depend on the index. The normalisation at infinity is that of \cite{PaperII} with $m$, since
$g-\operatorname{Log}z$ has no jump on $(-\infty,0)$ for every real $m$.

(ii) These are the jumps of the Hurwitz problem of Section~\ref{sec:hurwitz-thm} under the dictionary
$(n,\delta,\alpha+2\delta)\mapsto(m,-\gamma,\alpha-2\gamma)$, with the lattice band factor instead of the Hurwitz one and without the
moved nodes. The analytic estimates of that problem (the outer scalar, the a-priori bound, the comparison at the hard edge with
Bessel order $\alpha$ and arbitrary real Szeg\H{o} exponent, the band-factor gauge and the order-$K$ small-norm expansion, steps
(H-b)--(H-e) of \S\ref{sec:H-proof}) are statements for a fixed index in which the index enters only as a real scale, through
$t_m$, the Bessel variable $\zeta=m^2f$ and the moments $\mu_{m,j}$; none of them uses its integrality. The range of $\gamma$ is the
range $-1<\delta\le0$ of those estimates, with $|\tilde\alpha|<1$. The band factor transforms homogeneously:
$\mu_{n,j}(n/m)^{jp}z^j=\mu_{m,j}z^j$. The expansion then follows as in the proof of Theorem~\ref{thm:N6}(i), for the decoupled model at index $m$.
\end{proof}

The proof of Lemma~\ref{lem:frame}(ii) is a transfer of estimates, written out step by step against \cite{PaperII} and the
Hurwitz analysis in the notes of the repository named in the data availability statement; we do not repeat it. There is no
circularity: Section~\ref{sec:hurwitz-thm} uses Theorem~\ref{thm:shift} only for the natural form (Theorem~\ref{thm:H}(ii)), not
for the estimates.

\begin{proof}[Proof of Theorem~\ref{thm:shift}]
\emph{Two expansions of one number.} Fix $p\in(1,2)\setminus\mathbb Q$, $\alpha\in(-\frac12,\frac12)$, $\gamma\in(0,\frac14)$. As $n\to\infty$
through the integers, Theorem~\ref{thm:N6}(i) and Lemma~\ref{lem:frame}(ii), the latter multiplied by $(n/m)^p$, both express
$n^pb_n/c$ as a finite sum $\sum_ec_en^{-e}$ plus $o(n^{-E})$, for any given $E$ (choose $K,B,J$ large); in the $m$-frame the factors
$(n/m)^{(1+|\br|_w)p}$ and $m^{-k}$ are convergent series in $1/n$. In both, the exponents are $e=k+|\br|_wp$, $k\in\N_0$, and since $p$
is irrational, the exponents $k+jp$ are pairwise distinct. Two finite sums that agree up to $o(n^{-E})$ have the same
coefficients below $E$; at $e=k\in\N_0$ only $\br=0$ contributes. Hence, coefficientwise,
$B(n+\gamma;\alpha-2\gamma,\alpha)=(1+\gamma/n)^pB(n;\alpha,\alpha)$.

\emph{From an open set to an identity.} The coefficient of $n^{-k}$ in the difference of the two sides is a rational function
$R_k(p,\alpha,\gamma)$, which vanishes on $((1,2)\setminus\mathbb Q)\times(-\frac12,\frac12)\times(0,\frac14)$. Writing $R_k=N_k/D_k$ with
polynomials, $N_k(p,\cdot,\cdot)$ vanishes on an open set of $(\alpha,\gamma)$ for infinitely many $p$, so each of its coefficients, a
polynomial in $p$, has infinitely many zeros; thus $R_k\equiv0$. Renaming $\alpha\mapsto\nu$ and putting $\gamma=(\nu-\beta)/2$ gives
$B(n_0+\frac{\nu-\beta}2;\beta,\nu)=((n_0+\frac{\nu-\beta}2)/n_0)^pB(n_0;\nu,\nu)$ for every $\beta$. With
$F_p(N;\nu):=(N/n_0)^pB(n_0;\nu,\nu)$, $n_0:=N-\frac{\nu+1}2$, this is~(i).

\emph{The zeta layers.} The same comparison identifies the coefficient of $n^{-k-Jp}$ for each weight $J=|\br|_w$:
$\sum_{|\br|_w=J}\hat\mu^{\br}[X_{\br}]_k=0$, where $\hat\mu_j=2\zeta(-jp)/j$ and $X_{\br}:=(n/m)^{(1+J)p}\cD_{\br}(m;\alpha-2\gamma,\alpha)-\cD_{\br}(n;\alpha,\alpha)$.
For $J=1$ only $\br=e_1$ occurs, and $\zeta(-p)\ne0$ for $p\in(1,2)$, so $X_{e_1}=0$. For $J=2$ the relation reads
$\zeta(-2p)[X_{e_2}]_k+4\zeta(-p)^2[X_{2e_1}]_k=0$; the brackets are rational in $p$ (after removing a common factor $c^2$), and if
$[X_{e_2}]_k\not\equiv0$, then $\zeta(-2p)/\zeta(-p)^2$ would be a rational function of $p$. It is not, since it vanishes at every odd
integer $p\ge3$ ($\zeta(-2p)=0$, $\zeta(-p)\ne0$). So $X_{e_2}=X_{2e_1}=0$, and the identity argument extends both to all
$(p,\alpha,\gamma)$; (ii) follows as~(i) did.
\end{proof}

For weights $J\ge3$ the separation would need the linear independence of the monomials $\prod_j\zeta(-jp)^{r_j}$, $|\br|_w=J$, over
$\mathbb Q(p)$, which we do not prove; the shift of the layers of weight $\ge3$ is open. Exact computer algebra confirms the frame
identity~\eqref{eq:FI} identically in $(p,\alpha,\beta,\gamma)$ to order $n^{-4}$, and $40$-digit evaluation of the recursion to order
$n^{-8}$ at $32$ points, including points with $p<1$ and with $\gamma$ outside the range of Lemma~\ref{lem:frame}; controls with a
half shift $\beta=\alpha-\gamma$, the wrong sign, no frame factor, no index shift, or the Bessel order moved with $\beta$, all fail
(scripts \texttt{m4\_shift\_exact}, \texttt{sigma\_frame\_exact}, \texttt{sigma\_frame\_num}). These checks confirm the algebraic
consequence; they do not test the analysis of Lemma~\ref{lem:frame}(ii).

\subsection{The zeta layer as a Toda flow}\label{sec:toda}
In the unscaled variable the band factor $\exp(-\sum_k\mu_{n,k}z^k)$, $z=n^px$, is $\exp(\sum_kt_kx^k)$ with
\begin{equation}\label{eq:todatimes}
  t_k:=-\frac{2\zeta(-kp)}k ,
\end{equation}
a weight factor independent of $n$; these are the times of the Toda hierarchy. For a weight $w_t=e^{\sum_jt_jx^j}w$ with monic
orthogonal polynomials $\pi_n$, $h_n=\int\pi_n^2w_t$ and Jacobi matrix $\mathcal J$ (orthonormal basis), $\partial_{t_j}\pi_n$ has degree
$<n$, so $\partial_{t_j}h_n=(\mathcal J^j)_{nn}h_n$ and, since $b_n^2=h_n/h_{n-1}$,
\begin{equation}\label{eq:toda}
  \partial_{t_j}\log b_n^2=(\mathcal J^j)_{nn}-(\mathcal J^j)_{n-1,n-1},\qquad \partial_{t_1}a_n=b_{n+1}^2-b_n^2 .
\end{equation}

\begin{proposition}[Toda structure]\label{prop:toda}
Let $p\ne1$. The $\mu$-Taylor coefficients $\cD_{\br}$ of Theorem~\ref{thm:N6}(i), and their analogues for $a_n$, are the
Toda-hierarchy derivatives of the integer layers $b_n^{(0)}=cn^{-p}B(n)$ and $a_n^{(0)}=2cn^{-p}A(n)$, with $\partial_{\mu_j}=-n^{jp}\partial_{t_j}$.
For example, as identities of formal series in $1/n$ evaluated at $\mu=0$,
\begin{gather*}
  \partial_{\mu_1}(n^pb_n)=-\tfrac12n^p(a_n-a_{n-1})\,n^pb_n,\qquad
  \tfrac12\partial_{\mu_1}^2(n^pb_n)=\tfrac12n^{2p}\partial_{t_1}^2(n^pb_n),\\
  \partial_{\mu_2}(n^pb_n)=-\tfrac12n^{2p}\bigl[(\mathcal J^2)_{nn}-(\mathcal J^2)_{n-1,n-1}\bigr]n^pb_n,\qquad
  (\mathcal J^2)_{nn}=a_n^2+b_n^2+b_{n+1}^2,
\end{gather*}
with $\partial_{t_1}^2\log b_n^2=b_{n+1}^2-2b_n^2+b_{n-1}^2$.
\end{proposition}

\begin{proof}[Proof sketch]
Consider the tilted lattice measures $e^{\sum_j\tau_jx^j}\sigma_{p,s}$ with small real $\tau$. In the scaled variable the tilt is the
entire factor $\exp(\sum_j\tau_jn^{-jp}z^j)$, bounded on the contours of the frame; combined with the band factor it is again of the
gauge form of Lemma~\ref{lem:gauge}, with $\mu_{n,j}$ replaced by $\mu_{n,j}-\tau_jn^{-jp}$. So Theorem~\ref{thm:N6}(i) applies to
$b_n$ of the tilted measures uniformly for $\tau$ in a compact set, with layers polynomial in $\tau$, and the same holds for $a_n$
through the extraction $n^pa_n=(Q_2)_{12}/(Q_1)_{12}-(Q_1)_{22}$ of \cite{PaperII}. The hierarchy~\eqref{eq:toda} holds exactly for the
tilted measures at every $n$; inserting the expansions and taking finite differences in $\tau$ isolates each Taylor coefficient.
\end{proof}

The proof is only sketched: the $a_n$-extraction to every order and the uniformity in $\tau$ are not written out in detail. The
three displayed identities hold numerically to the central-difference error ($10^{-20}$ to $10^{-27}$) at three values of
$(p,\alpha)$ (script \texttt{m4\_toda}). In this reading the zeta layer is not new data: it is a fixed, $n$-independent Toda flow of the
integer layer, with times~\eqref{eq:todatimes} given by the zeta values. The diagonal integer layer behaves like $B$:

\begin{conjecture}[diagonal shift]\label{conj:diag}
For the decoupled model, the diagonal coefficient satisfies $n^pa(n;\beta,\nu)/(2c)=(n/M)^pG_p(M;\nu^2)$ with
$M:=n+\frac\beta2+1=N+\frac12$ and $G_p$ even in $1/M$ and independent of $\beta$, with
\[
  G_p(M;\nu^2)=1+\frac{a_2}{M^2}+O(M^{-4}),\qquad a_2=-\frac{p+1}{24}+\frac{(p+1)(p+2)\nu^2}8 .
\]
At $(p,\nu)=(2,-\frac12)$, $G_2=(1-1/(4M^2))^{-1}$.
\end{conjecture}

The evidence is numerical: the recursion at $40$ digits to order $8$, with $\beta$ and $\nu$ decoupled ($19$ checks), and $a_2$
reconstructed from $p\in\{\frac12,\dots,6\}$ and asserted at three further points (scripts \texttt{m4\_toda}, \texttt{m4\_diag}). The value at
the Bessel point is the diagonal of the Lommel measures. So $b_n$ is centred at $N$ and $a_n$ at $N+\frac12$, the arrangement of
\cite[Theorem~5.2]{BleherIts} for varying weights, and Theorem~\ref{thm:reflect} corresponds to the reversal symmetry
$b_n\leftrightarrow b_{-n}$, $a_n\leftrightarrow a_{-n-1}$ of the Toda lattice. The proof of Theorem~\ref{thm:shift} should apply to
$a_n$ through the same invariance of the extraction once the order-$K$ expansion of $a_n$ is written for $M(\beta,\nu)$.

\begin{remark}[the parity of the zeta layers]\label{rem:parity}
Given Proposition~\ref{prop:toda}, the parity statements of Conjecture~\ref{conj:S1Nprime} reduce to the integer layers. If $B$ is
even in $N$ and $A$ is even in $M=N+\frac12$, then $a_n-a_{n-1}=A(N+\frac12)-A(N-\frac12)$ (in natural form) is odd in $N$, so each
first-order zeta term is $\varepsilon_jN^{-1-jp}$ times an even series; second differences $b_{n+1}^2-2b_n^2+b_{n-1}^2$ are even, so the
$\varepsilon_1^2$ layer is $N^{-2-2p}$ times an even series, as found in Theorem~\ref{thm:N6}(ii). Higher Toda derivatives are built
from such differences and products, and we expect the full monomial structure to follow; we have not written this out. By
Corollary~\ref{cor:even} the evenness of $B$ is proved; that of $A$ is Conjecture~\ref{conj:diag}. The $1/(2j-1)$ in $g_j$ has a
heuristic reading in the same picture: for a slowly varying Jacobi matrix, $(\mathcal J^j)_{nn}$ is approximately
$\frac1{2\pi}\int_0^{2\pi}(a+2b\cos\theta)^j\diff\theta$, whose first variation off the point $a=2b$ involves
$\binom{2j-2}{j-1}/\binom{2j}j=j/(2(2j-1))$.
\end{remark}
